\section{Realización propietaria del estado enriquecido de Stein--Lyapunov}
\label{sec:ns-owner-closure}

La construcción de la \cref{sec:ns-pcf-prospective-construction} reúne en un
solo espacio de estado la torre PC--Fock completa, la memoria bidireccional,
el acarreo, la innovación microscópica, el canal espectral, la coordenada
homogénea de la fuerza, la unidad y la publicación excepcional. El tiempo
físico parametriza el almacenamiento terminal, mientras la profundidad
nonádica parametriza la recurrencia de Stein. La construcción que gobierna
esta sección es el \texttt{OWNER\_RECTOR} del estado enriquecido HMT. KYP,
NS--OBS y las cartas LRDN--LCN se conservan únicamente como
\texttt{CONTROL\_NO\_GOBERNANTE} de representaciones reducidas.

\subsection{Coligación PC--Fock y telescopía}

Sea
\[
 \mathbf x_{M,j}^{\rm own}
 \in\mathscr K_{M,j}^{\rm PCF}
 \widehat\otimes\mathcal E_{\rm exc}
 \oplus\mathscr G_M^{F,[j]}
\]
el estado homogéneo de la diagonal física. Para cada \(M,J,T\), la métrica
propietaria se construye prospectivamente como el Gram de la columna
terminal y de todas las pérdidas futuras de esa misma diagonal:
\begin{equation}
 \begin{aligned}
 \left\langle x,U_{M,j}^{\rm own}x\right\rangle
 :={}&
 \mathbb E_{j:J}
 \left\|\mathcal E_{M,J,T}^{\rm term}
       \Gamma_{M,j:J}^{\rm own}x\right\|^2\\
 &+\sum_{k=j}^{J-1}\mathbb E_{j:k}
 \left\|L_{M,k,T}^{\rm own}
       \Gamma_{M,j:k}^{\rm own}x\right\|^2 .
 \end{aligned}
 \label{eq:nsclose-owner-metric-definition}
\end{equation}
Todos los términos son cuadrados de las coordenadas ya construidas: grado
uno viscoso, grado dos convectivo, carry, innovación, archivo espectral,
memoria \(D/H\) y fuerza homogénea. Separar el primer paso de la suma da,
antes de introducir ninguna raíz KYP,
\begin{equation}
 \boxed{
 U_{M,j}^{\rm own}
 =
 (\Gamma_{M,j}^{\rm own})^*
 U_{M,j+1}^{\rm own}\Gamma_{M,j}^{\rm own}
 +(L_{M,j,T}^{\rm own})^*L_{M,j,T}^{\rm own},
 \qquad
 G_{M,J,T}^{\rm own}\preceq C_TU_{M,J}^{\rm own}.}
 \label{eq:nsclose-proprietary-terminal-gate}
\end{equation}
\label{eq:nsclose-owner-identity}
En la normalización terminal de \eqref{eq:ns-import-terminal-column} se puede
tomar \(C_T=1\). La igualdad de mapas, amplificada por la identidad en la capa
Moonshine, produce la cascada isométrica
\begin{equation}
 (\mathfrak U_{M,J,T}^{\rm own})^{\sharp}
 \mathfrak U_{M,J,T}^{\rm own}=I.
 \label{eq:nsclose-proprietary-full-cascade}
\end{equation}
Al iterar \eqref{eq:nsclose-owner-identity} se obtiene el telescopio
\begin{equation}
 U_{M,0}^{\rm own}
 =
 (\Gamma_{M,0:J}^{\rm own})^*U_{M,J}^{\rm own}
  \Gamma_{M,0:J}^{\rm own}
 +\sum_{j<J}(\Gamma_{M,0:j}^{\rm own})^*
  (L_{M,j,T}^{\rm own})^*L_{M,j,T}^{\rm own}
  \Gamma_{M,0:j}^{\rm own}.
 \label{eq:nsclose-owner-telescope}
\end{equation}
Su registro ortogonal de balance es
\begin{equation}
 \boxed{
 E_s(u_M(t))\le\|Z_{M,J,t}\|^2
 =\mathbb E_{0:J}G_{M,J,T}(t)
 +\sum_{j<J}\mathbb E_{0:j}
   \Lambda_{M,j}^{\rm own}([0,t]).}
 \label{eq:nsclose-proprietary-ledger}
\end{equation}
La raíz común
\(\Xi_{M,0}=J_{0\to M}\Xi_0\), la conservación evolutiva de la unidad y el
puerto de carry de \eqref{eq:ns-import-carry-port} hacen que el lado derecho
se pague una sola vez. El presupuesto propietario es
\begin{equation}
 \begin{aligned}
 \mathcal B_T^{\rm own}:={}&
 C_{\rm HMT}
 +(1+2\widetilde M_{s,\beta,\nu}^2)
  e^{R_{\rm F}^2\|u_0\|_{H^s}^2}\\
 &+C_{\nu,s}\int_0^T\|f(t)\|_{H^{s-1}}^2\,dt
 +\|\mathcal J_T^{\rm own}(u_0,f)\|^2 ,
 \end{aligned}
 \label{eq:nsclose-owner-data-budget}
\end{equation}
donde \(\mathcal J_T^{\rm own}\) es la inyección de raíz determinada por el
estado inicial transportado por TPK, sus dos hojas y la entrada forzada. Sus
amplificaciones son
\[
 \mathcal J_{M,J,T}^{\rm own}
 :=I_{0\to M,J}\mathcal J_T^{\rm own},
 \qquad I_{0\to M,J}^*I_{0\to M,J}=I.
\]
Por \eqref{eq:nsclose-owner-telescope}, no es una cantidad reconstruida desde
la órbita futura. La conservación de la unidad y la coisometría del lector
two-way dan, sobre el dominio declarado,
\begin{equation}
 \sup_{M,J}\|\mathcal J_{M,J,T}^{\rm own}(u_0,f)\|^2
 \le C_{{\rm HMT},T}
 \left(1+\|u_0\|_{H^{s+1}}^2
 +\int_0^T\|f(t)\|_{H^{s-1}}^2\,dt\right).
 \label{eq:nsclose-owner-root-uniformity}
\end{equation}
Así \(\mathcal B_T^{\rm own}\) es independiente de \(M,J\). El ledger da
\begin{equation}
 \boxed{
 \sup_{M,J}\left[
  \sup_{0\le t\le T}\|Z_{M,J,t}\|^2
  +\frac\nu2\int_0^TD_{s,M}(t)\,dt
 \right]\le\mathcal B_T^{\rm own}<\infty.}
 \label{eq:nsclose-proprietary-uniform-budget}
\end{equation}

La publicación hidrodinámica es el retorno de primer grado de la misma
cascada:
\begin{equation}
 R_{M,J,T}^{\rm full}
 :=\mathcal R_{1,M}
   (\mathfrak U_{M,J,T}^{\rm own})^{\sharp},
 \qquad
 \boxed{
 R_{M,J,T}^{\rm full}Z_{M,J,t}=u_M(t).}
 \label{eq:nsclose-proprietary-exact-return}
\end{equation}
\label{eq:nsclose-exact-full-return}
Esta igualdad coincide con \eqref{eq:ns-import-exact-return} y evita la
evaluación de una rama individual del refinamiento.

\begin{theorem}[Estimación uniforme de la realización propietaria PC--Fock]
\label{thm:nsclose-proprietary-full}
Para todo \(T<\infty\), las aproximaciones de Galerkin satisfacen
\begin{equation}
 \sup_M\sup_{0\le t\le T}\|u_M(t)\|_{H^s}^2
 +\nu\sup_M\int_0^T\|u_M(t)\|_{H^{s+1}}^2\,dt
 \le C_s\mathcal B_T^{\rm own},
 \label{eq:nsclose-uniform-sobolev-precompactness}
\end{equation}
\label{eq:nsclose-uniform-sobolev}
con una constante independiente de \(M\) y de la profundidad nonádica.
\end{theorem}

\begin{proof}
La identidad \eqref{eq:nsclose-proprietary-exact-return} y la contractividad
del retorno dan
\(\|u_M(t)\|_{H^s}^2\le C_s\|Z_{M,J,t}\|^2\). La equivalencia de
\(D_{s,M}\) con la norma \(H^{s+1}\) en el subespacio solenoidal de media
cero, junto con \eqref{eq:nsclose-proprietary-uniform-budget}, prueba la
estimación.
\end{proof}

\subsection{Control alternativo por métrica de ruta}

La carta que sigue comprueba el refinamiento en coordenadas normalizadas. No
construye ni condiciona \(U^{\rm own}\): omite la unidad, la publicación
excepcional y parte del canal afín que ya residen en
\eqref{eq:nsclose-owner-metric-definition}.

La carta de Ambivalencia aporta
\begin{equation}
 G_{\rm av}=\begin{pmatrix}2&-1\\-1&1\end{pmatrix},
 \qquad
 M_{\rm av}=\begin{pmatrix}3&-1\\1&0\end{pmatrix},
 \qquad
 M_{\rm av}^*G_{\rm av}M_{\rm av}
 \preceq\varphi_{\rm HMT}^4G_{\rm av}.
 \label{eq:nsclose-ambivalence-route-metric}
\end{equation}
Después de la normalización \(\widehat M_j=M_j/3\),
\begin{equation}
 \mathsf R_j:=G_j-\widehat M_j^*G_{j+1}\widehat M_j\succeq0.
 \label{eq:nsclose-route-exact-defect}
\end{equation}
En la carta estacionaria,
\[
 \mathsf R_{\rm av}
 =\frac19\begin{pmatrix}5&-4\\-4&7\end{pmatrix},
 \qquad
 \lambda_{\min}(\mathsf R_{\rm av})
 =\frac{6-\sqrt{17}}9.
\]
Con la métrica de grafo
\[
 \mathsf G_M^{\rm gr}
 =I+(\mathcal Y_{\beta,M})^*\mathcal Y_{\beta,M},
 \qquad
 \mathcal C_M
 =\mathcal Y_{\beta,M}(\mathsf G_M^{\rm gr})^{-1/2},
\]
definimos
\begin{equation}
 \mathcal A_{M,j}=\widehat M_j\otimes I,
 \qquad
 \mathcal O_{M,j}=\mathsf R_j^{1/2}\otimes\mathcal C_M,
 \qquad
 \mathcal D_{M,j}
 =\mathsf R_j^{1/2}\otimes(I-\mathcal C_M^*\mathcal C_M)^{1/2}.
 \label{eq:nsclose-owner-route-column}
\end{equation}
La identidad de Stein de ruta es
\begin{equation}
 \boxed{
 \mathsf U_{M,j}^{\rm PCF}
 -\mathcal A_{M,j}^*\mathsf U_{M,j+1}^{\rm PCF}\mathcal A_{M,j}
 =\mathcal O_{M,j}^*\mathcal O_{M,j}
  +\mathcal D_{M,j}^*\mathcal D_{M,j}.}
 \label{eq:nsclose-owner-stein}
\end{equation}
Su iteración da
\begin{equation}
 \begin{aligned}
 \mathsf U_{M,0}^{\rm PCF}
 -\mathcal A_{M,0:J}^*\mathsf U_{M,J}^{\rm PCF}\mathcal A_{M,0:J}
 =\sum_{j<J}\mathcal A_{M,0:j}^*
 \bigl(\mathcal O_{M,j}^*\mathcal O_{M,j}
      +\mathcal D_{M,j}^*\mathcal D_{M,j}\bigr)
 \mathcal A_{M,0:j}.
 \end{aligned}
 \label{eq:nsclose-owner-stein-telescope}
\end{equation}
Esta factorización describe la geometría del refinamiento. La recurrencia
afín \eqref{eq:ns-import-affine-stein} incorpora además la evolución física,
los términos cruzados de Carleman y la fuerza.

La memoria bidireccional convierte la pérdida de ruta en un presupuesto. Con
\(\mathcal N_{M,j}=2(r-1)M_{M,j}^-\),
\begin{equation}
 E_j\Delta_T\mathcal N_{M,j+1}
 =\Delta_T\mathcal N_{M,j}+2K_{M,j}(T),
 \qquad
 K_{M,j}(T)=\int_0^T\kappa_{M,j}(t)\,dt.
 \label{eq:nsclose-carry-memory-telescope}
\end{equation}
La telescopía y la cota de la raíz producen
\begin{equation}
 \mathbb E_{0:J}G_{M,J,T}(0)
 \le E_{s,M}(P_Mu_0)+|\Delta_T\mathcal N_0|
 +\|\Xi_{M,0}\|^2+Q_f(T),
 \label{eq:nsclose-owner-terminal-bound}
\end{equation}
que es la forma abreviada del control de ruta
\eqref{eq:ns-import-root-bound}--\eqref{eq:ns-import-explicit-budget}. Esta
carta registra el coste de la hoja negativa,
\[
 \Delta_T\mathcal N_{M,0}
 =2\int_0^T B^-_{s,M,0}(u_M(t))\,dt.
\]
En esta representación reducida, la desigualdad
\[
 X^-_{M,T}(0)\preceq C_T\mathsf U^{\rm PCF}_{M,0},
 \qquad \sup_M C_T<\infty
 \tag{NS--OBS}
\]
es un falsador útil de la carta. Su fallo no se transporta al owner, porque
\(\Delta_T\mathcal N_{M,0}\) no sustituye a la inyección de raíz
\(\mathcal J_T^{\rm own}\) de
\eqref{eq:nsclose-owner-data-budget}--\eqref{eq:nsclose-owner-root-uniformity}.
Análogamente, KYP y LRDN--LCN verifican completaciones o dominaciones de
lectores reducidos; ninguna de ellas es una premisa del
\cref{thm:nsclose-proprietary-full}.

\subsection{Existencia, unicidad y suavidad global}

\begin{theorem}[Regularidad global de Navier--Stokes]
\label{thm:nsclose-global-smooth}
\label{thm:realizaciones-ns}
Sea \(s>5/2\), \(\nu>0\), sea
\(u_0\in H^{s+1}_\sigma(\mathbb T^3)\) de media cero y sea
\(f\in L^2_{\rm loc}([0,\infty);H^{s-1}_\sigma)\). Existe una única solución
global
\begin{equation}
 u\in L^\infty_{\rm loc}([0,\infty);H^s_\sigma)
 \cap L^2_{\rm loc}([0,\infty);H^{s+1}_\sigma),
 \qquad
 \partial_tu\in L^2_{\rm loc}([0,\infty);H^{s-1}_\sigma).
 \label{eq:nsclose-global-regularity-class}
\end{equation}
Para dato y fuerza suaves, \(u\) y la presión normalizada son suaves.
\end{theorem}

\begin{proof}
El \cref{thm:nsclose-proprietary-full} da acotación uniforme de
\((u_M)_M\) en \(L^\infty(0,T;H^s)\cap L^2(0,T;H^{s+1})\). La ecuación de
Galerkin y la continuidad
\[
 H^s\times H^s\longrightarrow H^{s-1},
 \qquad (u,v)\longmapsto\mathbb P(u\cdot\nabla v),
\]
acotan \(\partial_tu_M\) en \(L^2(0,T;H^{s-1})\). Como
\(H^{s+1}\Subset H^s\hookrightarrow H^{s-1}\), Aubin--Lions produce, tras
extraer una subsucesión,
\[
 u_M\longrightarrow u\quad\hbox{fuertemente en }L^2(0,T;H^s).
\]
Esta convergencia permite pasar al límite en
\(\mathbb P(u_M\cdot\nabla u_M)\), mientras la convergencia débil conserva el
término viscoso y el dato inicial. Se obtiene una solución fuerte en
\([0,T]\).

Si \(u,v\) son dos soluciones y \(w=u-v\), la estimación de producto y
Young dan
\[
 \frac12\frac d{dt}\|w\|_{H^{s-1}}^2
 +\frac\nu2\|w\|_{H^s}^2
 \le C_{s,\nu}\bigl(\|u\|_{H^s}^2+\|v\|_{H^s}^2\bigr)
 \|w\|_{H^{s-1}}^2.
\]
Grönwall prueba unicidad. Al repetir la construcción en los horizontes
\(T=1,2,\ldots\), las soluciones coinciden en los intervalos solapados y se
obtiene la solución global.

La presión se recupera de
\begin{equation}
 -\Delta p=\partial_i\partial_j(u_i u_j)-\operatorname{div}f,
 \qquad \int_{\mathbb T^3}p\,dx=0.
 \label{eq:nsclose-pressure-recovery}
\end{equation}
La iteración de la estimación en \(s+k\) y la ecuación recuperan todas las
derivadas espaciales y temporales cuando \(u_0\) y \(f\) son suaves.
\end{proof}

La contribución estructural consiste en representar la convección como una
salida lineal de la torre PC--Fock, conservarla en la memoria \(D/H\) e
integrarla en una identidad ortogonal con retorno físico exacto. La
estimación uniforme resultante se publica después como regularidad global de
la ecuación hidrodinámica.
